\section{Related Work}
\subsection{Privacy and trustworthy crowdsensing in vehicular networks}
Research on vehicular crowdsensing security has addressed participant identity, confidentiality of sensed data, and evidence trustworthiness. AARE combines anonymous authentication with privacy-preserving assessment of data reliability~\cite{Cheng2023}. BARM introduces blockchain-assisted anonymous authentication and reputation management to evaluate and update participant reputations~\cite{Zheng2025BARM}. Regarding data processing, PPVDA studies privacy-preserving and verifiable aggregation~\cite{Zhou2024}; BCS-LPP employs techniques including Geohash and order-preserving encryption to protect sensing locations while supporting location checks and task management~\cite{Shen2024}. Cooperative detection of anomalous vehicular behavior has also been investigated at both vehicle and roadside-unit levels~\cite{Sultana2024}. Other recent work examines lightweight anonymous authentication, certificateless aggregate signatures, and SNARK-based verification of computations on authenticated data~\cite{Imam2025JISA,Wang2025JISA,Reinhart2026ADSC}. Recent work in \emph{Cybersecurity} extends this authentication perspective in different settings: Umrani and Palmieri investigate threshold password-authenticated key exchange for vehicular ad hoc networks~\cite{CybersecUmrani2025}, while Wang et al. develop privacy-preserving multi-authority aggregate signatures for smart-grid message authentication~\cite{CybersecWang2026}. These systems address secure communication or message authentication, rather than the physical coverage of uncertain evidence.

These approaches strengthen privacy and trust at different points in the sensing workflow, but they certify different objects. Identity authentication establishes participation rights; private aggregation protects inputs and checks computation; and misbehavior detection identifies suspicious participants. In the setting addressed here, several authenticated sources provide uncertain or conflicting spacetime observations about one event. Source legitimacy and signature validity alone do not establish whether the event complies with a given policy. The additional requirement is a verifiable basis for a policy decision that retains the uncertainty of the individual observations.

\subsection{Zero-knowledge location and geometric-policy verification}
Location privacy can be pursued through application-specific obfuscation and protected queries~\cite{Mdee2022,Atmaca2024}, or by credentials and authentication protocols that limit the disclosure of sensitive location information~\cite{Ghani2025,Zhang2025Sybil,Wang2025ZKAuth}. Zero-knowledge proofs offer another approach: a verifier can check a geometric constraint without learning the protected coordinate or an entire trajectory. ZKLP, proposed by Ernstberger and colleagues, incorporates floating-point geometry into zero-knowledge proofs for region membership and proximity predicates~\cite{ZKLP}. The vehicle-policy protocols of Bogdanov and colleagues instead use locations signed by a trusted witness to prove compliance with driving-area requirements in subsidy or taxation applications~\cite{VehiclePolicy}. Privacy-preserving policy checks have also been studied beyond location services. Kerl et al. combine attribute-based access control with privacy-preserving evaluation of protected attributes~\cite{CybersecKerl2025}, whereas Zhang applies zero-knowledge proofs to policy enforcement on encrypted network traffic~\cite{CybersecZhang2026}. Their verified statements concern attributes or traffic rules rather than threshold-supported uncertainty across signed spacetime observations.

These studies demonstrate that geometric predicates and vehicle-trajectory policies can be certified while withholding the underlying positions. Protecting a position, however, differs from establishing that the asserted position is adequately supported by observations. With disagreeing measurement intervals or sources that may maliciously sign data, proving membership of one chosen point or an existing trajectory does not establish that \emph{all} states supported by sufficient evidence satisfy the policy. Our proof relation therefore concerns the threshold support induced jointly by multiple signed intervals, rather than repeated membership testing of a selected coordinate.

\subsection{Byzantine-robust sensor fusion}
Uncertainty and disagreement across measurements have long been studied in fault-tolerant sensor fusion. Marzullo's approach to continuous-valued sensor failures uses overlapping intervals to characterize states consistent with potentially faulty observations~\cite{Marzullo1990}. BFT-PoLoc applies Byzantine-resistant geometry to localization from Internet delays~\cite{BFTPoLoc}. PG combines private computation and Byzantine-tolerant sensor fusion while addressing output delivery~\cite{PG}. CRSF further examines private fusion when sensors and servers may collude, including participation consistency and protocol liveness~\cite{CRSF}.

Unlike protocols that verify an individual location statement, these systems address the difficulties posed by inconsistent observations and malicious participants. Some offer stronger participation or delivery guarantees than a fixed issuer roster. Their main output, however, is a fused estimate or a secure computation result, rather than a noninteractive proof tied to a particular report that an external verifier can use to check spacetime-policy compliance. We draw on fault-tolerant modeling but ask a different question: which states reach a predefined source-support threshold, and how can their relationship with a public policy be certified without disclosing the underlying signed intervals?

\subsection{Comparison and research positioning}
Table~\ref{tab:related-work} compares representative methods by the object checked or computed, the evidence and trust model, and the privacy and output interface. As the methods use different adversarial models and assumptions, the comparison is intended to clarify their roles, not to rank their security or performance.

\begin{table}[htbp]\centering\footnotesize
\caption{Functional scope and trust assumptions of representative systems.}\label{tab:related-work}
\begin{tabularx}{\linewidth}{@{}p{0.16\linewidth}p{0.23\linewidth}p{0.29\linewidth}X@{}}
\toprule Method & Object & Evidence and trust assumptions & Privacy and output\\\midrule
BARM~\cite{Zheng2025BARM} (2025) & Identity and reputation & Reputation evaluation and blockchain-assisted identity management & Anonymous authentication and private reputation evaluation\\
PPVDA~\cite{Zhou2024} (2024) & Sensing aggregates & Homomorphic message authentication, secret sharing, malicious participants & Private inputs and verifiable aggregation\\
BCS-LPP~\cite{Shen2024} (2024) & Sensing location & Geohash, order-preserving encryption, deposit mechanism & Location privacy and task/location checking\\
BFT-PoLoc~\cite{BFTPoLoc} (2024) & Delay-based location & Internet delays and Byzantine-resistant geometry & Ranging and multilateration; privacy not rated here\\
PG~\cite{PG} (2024) & Fused sensor result & Malicious sensors and aggregating servers & Private inputs/results; output-delivery guarantee\\
CRSF~\cite{CRSF} (2026) & Fusion and participation & Sensor--server collusion and consistency of submissions & Private fusion, explicit abort and liveness\\
ZKLP~\cite{ZKLP} (2025) & Region and proximity & Hidden coordinates; floating-point geometry & Zero-knowledge geometric proof\\
Vehicle policy~\cite{VehiclePolicy} (2025) & Driving-area policy & Locations signed by a trusted witness & Hidden trajectory; subsidy and tax compliance\\
\textbf{This work} & \textbf{All threshold-supported spacetime states} & \textbf{Signed uncertainty boxes, fixed roster, budgets $f$ and $r$} & \textbf{Hidden boxes/signatures; report-bound four-valued proof}\\\bottomrule
\end{tabularx}
\par\smallskip\raggedright\footnotesize BFT-PoLoc and CRSF refer to the cited arXiv versions. All privacy and correctness claims depend on their respective trust assumptions. Semantic inconsistency must not be confused with abort on missing evidence or proof failure.
\end{table}

Zero-knowledge location verification and Byzantine-robust fusion thus provide complementary foundations. The former can certify geometric constraints over undisclosed positions, but generally does not take threshold-supported collections of signed uncertainty intervals as the proof object. The latter reasons about anomalies and uncertain inputs, but does not primarily aim to issue a third-party-verifiable policy-compliance certificate. We connect these directions: threshold support determines the states permitted by the uncertain evidence, and zero knowledge certifies their relationship with the public policy.

Compared with point- or trajectory-based proofs, this design adds issuer-signature checks and set-level compliance constraints, at the cost of larger circuits and a pre-established issuer roster. Compared with private sensor fusion, it provides a noninteractive certificate bound to the task, report, and challenge. This certificate remains conditional on trusted enrollment, the fault budget, and physical coverage; the fixed roster does not guarantee dynamic participation or issuer availability. The next sections make these premises and their implementation costs explicit.
